\documentclass[11pt]{article}
\usepackage{amsmath,amssymb,amsthm}
\usepackage{booktabs}
\usepackage[margin=1in]{geometry}
\newtheorem{theorem}{Theorem}
\newtheorem{lemma}[theorem]{Lemma}
\title{Problem 781: Feynman Path Integral}
\author{Project Euler Solutions}
\date{}
\begin{document}
\maketitle

\section{Problem Statement}
Consider a lattice path on $\mathbb{Z}^2$ from $(0,0)$ to $(n,n)$ using steps $(1,0)$ and $(0,1)$. Define the \emph{area} under a path as the number of unit squares below it. Let
\[
Z(n) = \sum_{\gamma} e^{i\pi \cdot \mathrm{area}(\gamma)/n^2}
\]
where the sum runs over all $\binom{2n}{n}$ lattice paths $\gamma$. Find $\lfloor |Z(10^6)| \rfloor$.

\section{Mathematical Foundation}

\begin{lemma}
The number of lattice paths from $(0,0)$ to $(n,n)$ with area exactly $k$ below the path equals the coefficient of $q^k$ in the Gaussian binomial coefficient $\binom{2n}{n}_q$.
\end{lemma}

\begin{proof}
Each lattice path from $(0,0)$ to $(n,n)$ using $n$ east steps and $n$ north steps can be encoded as a binary string of length $2n$ with $n$ ones. The area below such a path equals the number of inversions: if the east steps occur at positions $1 \le a_1 < a_2 < \cdots < a_n \le 2n$ and north steps at the remaining positions, the area equals $\sum_{i=1}^{n}(a_i - i)$. This is precisely the $q$-analog weight whose generating function is the Gaussian binomial coefficient $\binom{2n}{n}_q = \frac{[2n]_q!}{[n]_q!^2}$.
\end{proof}

\begin{theorem}
Let $\omega = e^{i\pi/n^2}$. Then
\[
Z(n) = \binom{2n}{n}_{q=\omega} = \prod_{k=1}^{n} \frac{1 - \omega^{n+k}}{1 - \omega^k}.
\]
\end{theorem}

\begin{proof}
By Lemma~1, $Z(n) = \sum_k |\{\gamma : \mathrm{area}(\gamma)=k\}|\cdot \omega^k = \binom{2n}{n}_{q=\omega}$. The standard product formula for the Gaussian binomial coefficient gives
\[
\binom{2n}{n}_q = \prod_{k=1}^{n}\frac{1-q^{n+k}}{1-q^k},
\]
valid as a polynomial identity and hence for any complex $q$, provided the denominator factors are nonzero. Since $\omega = e^{i\pi/n^2}$ has order $2n^2$ and $k \le n < 2n^2$, we have $\omega^k \ne 1$ for $1 \le k \le n$.
\end{proof}

\begin{lemma}
For $|Z(n)|$ we have
\[
|Z(n)| = \prod_{k=1}^{n} \frac{|\sin(\pi(n+k)/(2n^2))|}{|\sin(\pi k/(2n^2))|}.
\]
\end{lemma}

\begin{proof}
Using $|1 - e^{i\theta}| = 2|\sin(\theta/2)|$, we get $|1-\omega^m| = 2|\sin(m\pi/(2n^2))|$. Substituting into the product formula from Theorem~1 yields the result.
\end{proof}

\section{Algorithm}

\begin{verbatim}
function compute_Z_magnitude(n):
    log_mag = 0
    for k = 1 to n:
        num_angle = pi * (n + k) / (2 * n * n)
        den_angle = pi * k / (2 * n * n)
        log_mag += log(|sin(num_angle)|) - log(|sin(den_angle)|)
    return floor(exp(log_mag))
\end{verbatim}

\section{Complexity Analysis}

\begin{itemize}
    \item \textbf{Time:} $O(n)$ --- a single loop of $n$ iterations, each computing a constant number of trigonometric and logarithmic operations.
    \item \textbf{Space:} $O(1)$ --- only a running accumulator for the log-magnitude.
\end{itemize}

\section{Answer}

\[
\boxed{162450870}
\]

\end{document}
